\section*{Acknowledgments}

The author thanks the institutions whose commitment to open public data made this work possible: the World Bank, the International Labour Organization, the International Monetary Fund, the OECD, UNU-WIDER, the United Nations Population Division, the Penn World Table project, and the national statistical institutes of Peru, Chile, Colombia, and Mexico.

This paper is the first concrete step of a longer research path, and it owes much to the author's family: to his sister, Luciana Franceska Scavino Alfaro, for walking beside him along the way; to his mother, Isabel Betzabe Alfaro Ninahuanca, for the drive to reach it; and to his father, Luigi Stefano Scavino Montero, for the light that showed the way. He also thanks Oddie for the steady company.

\section*{Data and code availability}

The primary specification, including calibrated parameter supports, was fixed before the baseline run (frozen specification; SHA-256 below).\par
\nolinkurl{0080d502a50db430998b56ebc6419ee181c90a1331cd2b650bb25db298523f81}

The horizon-inversion statistic $H^*$ was introduced after the baseline as a derived reporting measure; it changes no calibrated parameter and no baseline classification.

The repository contains the complete replication code and archived inputs with SHA-256 manifests. Reproduction instructions appear in \nolinkurl{REPRODUCING.md}, while \nolinkurl{DATA_AVAILABILITY.md} records source-specific licenses and redistribution decisions. Full technical documentation, parameters, and diagnostics are provided in the replication package, archived at \url{https://doi.org/10.5281/zenodo.21348534}.

\section*{Declaration of competing interest}

The author declares no competing financial interests or personal relationships that could have appeared to influence the work reported in this paper.

\section*{Funding}

This research received no external funding.

\section*{Declaration of Generative AI and AI-assisted technologies in the writing process}

During the preparation of this work the author used Claude (Anthropic) in order to improve the readability, language, and formatting of the manuscript. After using this tool, the author reviewed and edited the content as needed and takes full responsibility for the content of the publication.
